%% ===================================================================
\section{Discussion}
\label{sec:discussion}

\paragraph{The tile exponent is not an environmental constant}
The central premise of descriptor-based tile transfer is that a tile has a
path-loss exponent determined by its environment, which can be read from
descriptors and carried to similar tiles elsewhere. On this corpus the
premise is weak, for three measured reasons. Two estimates of the same tile's
exponent from disjoint halves of the sites differ by \seSplitDiff{} in
standard deviation, so the exponent depends on which sites serve the tile;
it depends on the dynamic range of the measurements, shifting by
\censDgOneForty{} when the cut-off is lowered by \SI{10}{\decibel}; and the
cross-city spread of the city medians (\betweenSd) is small against the
spread within a city. A tile is seen by each site through a different street
canyon, and the exponent fitted over all of them mixes the geometry of the
tile with that of the paths reaching it. This explains why learned operators
improve on the source median only modestly, why few-shot correction of the
exponent needs tens of surveyed tiles before it helps, and why a planning
tool fed with even the \emph{true} receiver-tile labels assigns half of the pixels to
the wrong site and draws cell edges \SI{\edgeOracle}{\metre} away from the
true ones.

\paragraph{What to transfer}
The quantity that transfers is the effect of the geometry along a link,
not a per-tile summary of it. Site-aware link transfer uses the same
descriptors, the same training cities and the same protocol, but conditions
on the path between site and receiver; it improves path loss, serving site,
coverage and planning in every held-out city and outperforms the tile model
with perfect labels, however the tile parameters are composed along the
path (Section~\ref{sec:res-mixed}). It keeps its advantage over transferred tile parameters
at another frequency, with other sites and on independently modelled scenes.
On the latter, whose materials differ from the training corpus, it is biased
until a survey of 25 random tiles calibrates it, after which its path loss
matches exact tile labels and its serving-site accuracy exceeds theirs: the
material models of the training corpus limit what transfers zero-shot, and a
small unbiased survey corrects the rest. Its most important features, the knife-edge loss and
the height of the obstruction above the line of sight, are the quantities of
classical diffraction theory~\cite{itup526}, so the model remains
interpretable in physical terms. Two design choices matter for planning.
First, the model must represent the absence of a link: the hurdle, which
predicts whether a site is received at all, reduces the planning regret from
\regretLinkReg\% to \regretLink\% at the same path-loss error. Second, its
errors must be stated at the level of the decision: the city-level
certificate turns the leave-one-city-out residuals into a coverage promise
with a guaranteed risk, and the more accurate the model, the more of the
city it can promise (\certClaimLink\% against \certClaimGbdt\% for tile
transfer). For a digital twin this means that the transferable
representation should be attached to links (or to site--tile pairs), which
is also what a planning tool evaluates.

\paragraph{Kernel transfer}
Theorem~\ref{thm:regimes} identifies two limits of kernel transfer; the
tuned operator is in neither, but with the bandwidths selected on two
targets it approaches the source mean without reaching the flat regime.
Its disadvantage against the learned regressors is the noise-transmission
term of Proposition~\ref{prop:decomp} combined with an isotropic distance:
learning one length scale per descriptor (GP-ARD) removes most of the gap,
which locates the weakness in the metric rather than in kernel averaging as
such.

\paragraph{Implications for planning}
Four practical conclusions follow. First, a transfer model for a
digital twin should be validated at the network level; parameter errors of
similar size lead to very different planning losses (Table~\ref{tab:network}),
and a uniformly biased model such as the 3GPP UMa NLoS constants has a
path-loss error of \SI{\netUmaNlos}{\decibel}, far above all transfer
operators, yet a planning
regret of \regretUmaNlos\%, below that of transferred tile parameters,
because a uniform bias does not change which sites cover the most pixels. Second, models trained on links below a
receiver cut-off overstate the reach of distant sites unless the cut-off is
modelled; every zero-shot tile model with transferred parameters predicts that a
single site covers \covTarget\% of a city. Third, small surveys do not help a model without systematic error, but
correct a shifted one, and they must be unbiased: a design that maximises
the variance reduction of the offsets selects unrepresentative tiles. Fourth, coverage claims should come with a margin estimated across
cities: without it, \certNaiveFalseGbdt\% of the pixels promised by tile
transfer are not covered.

\paragraph{Limitations}
The path loss is ray traced rather than measured, so absolute values depend
on the solver, the ITU material models and the level-of-detail-1 geometry;
the results are comparisons under identical conditions. The external scenes
test a different geometry source and different materials but the same
solver, and the siting replicate and the second band test other sites and
another frequency; none replaces validation against measurements, which the
public drive-test data we could obtain do not allow because they lack site
positions. Ray tracing is nevertheless a demanding rather than a lenient
testbed for tile transfer: its losses arise site by site from the building
geometry of each path, not from a statistical model with tile structure, so
a tile model has no advantage built in, and it provides what measurements
cannot, a complete map for every site and exact labels for every tile, which
separate the error of the tile model from that of transfer. Between \impMin\% and \impMax\% of the building heights are
imputed per city. With eleven cities, city-level tests have limited power:
an exact Wilcoxon test with Holm's correction over a dozen comparisons
requires nearly every city to agree, which is why no improvement of a tile
operator over the median survives it although the learned operators improve
on it in eight or nine of eleven cities. The corpus has flat terrain, \nSitesPerCity{} rooftop sites
per city with isotropic antennas, full load and a fixed link budget. Tile
labels come from the same ray tracer as the network truth, so the oracle is
the reference for tile models with exact labels on this corpus, not a
statement about measured data.
